\documentclass[sigconf]{acmart}

\usepackage{graphicx}
\usepackage{booktabs}
\usepackage{hyperref}
\usepackage{url}
\usepackage{caption}
\usepackage{subcaption}
\usepackage{enumitem}
\usepackage{multirow}
\usepackage{array}
\usepackage{tabularx}

\AtBeginDocument{%
  \providecommand\BibTeX{{%
    \normalfont B\kern-0.5em{\scshape i\kern-0.25em b}\kern-0.8em\TeX}}}

% Copyright, conference, and ISBN details removed — to be added before submission
\setcopyright{none}
\renewcommand\footnotetextcopyrightpermission[1]{} % suppress copyright footnote

% Suppress running headers (venue/date line)
\acmConference{}{}{}
\acmBooktitle{}
\acmPrice{}
\acmISBN{}
\acmDOI{}

% Suppress ACM Reference Format block and running headers
\settopmatter{printacmref=false, printfolios=true}

\begin{document}

\title{LLM-Based Intelligent Notification Composition: From Static Personalization to Context-Aware Persuasive Messaging}

\author{Nilesh Agrawal}
\email{nilesh.d.agrawal@gmail.com}
\orcid{0000-0000-0000-0000}
\affiliation{%
  \institution{Independent Researcher}
  \city{Seattle}
  \state{WA}
  \country{USA}
}

\renewcommand{\shortauthors}{}
\fancyhead{} % clear all header fields
\fancyhead[L]{\textit{LLM-Based Intelligent Notification Composition}}
\fancyhead[R]{\textit{Agrawal, 2026}}

\begin{abstract}
Modern notification stacks have invested heavily in three of four key decisions---audience selection, content ranking, and send-time optimization---but the fourth decision, \textit{how to communicate the recommendation}, remains the least optimized stage. This paper argues that message quality is an independent, underinvested lever and that Large Language Models (LLMs) provide a principled mechanism for improving it. We make three contributions absent from prior work: (1) a six-dimension Message Quality Framework with empirical grounding; (2) an Architectural Attribution model that disentangles message generation from adjacent recommender components; and (3) a Binding-Constraint Framework specifying when LLM generation is the binding constraint and when templates remain superior. Synthesizing 28 sources screened from 142 records across academic and industrial deployments, we find that CTR improvements from LLM-based message generation range from +1\% to +14.5\% over static templates, with the variance driven primarily by baseline sophistication. Improvements concentrate along the freshness, contextual relevance, and actionability dimensions. The binding-constraint framework offers practitioners a principled basis for making deployment decisions, including when not to use LLMs at all.
\end{abstract}

\begin{CCSXML}
<ccs2012>
   <concept>
       <concept_id>10002951.10003317.10003347.10003350</concept_id>
       <concept_desc>Information systems~Recommender systems</concept_desc>
       <concept_significance>500</concept_significance>
       </concept>
   <concept>
       <concept_id>10010147.10010178.10010179.10010182</concept_id>
       <concept_desc>Computing methodologies~Natural language generation</concept_desc>
       <concept_significance>500</concept_significance>
       </concept>
   <concept>
       <concept_id>10003120.10003130.10011762</concept_id>
       <concept_desc>Human-centered computing~Empirical studies in collaborative and social computing</concept_desc>
       <concept_significance>300</concept_significance>
       </concept>
 </ccs2012>
\end{CCSXML}

\ccsdesc[500]{Information systems~Recommender systems}
\ccsdesc[500]{Computing methodologies~Natural language generation}
\ccsdesc[300]{Human-centered computing~Empirical studies in collaborative and social computing}

\keywords{Large Language Models, Push Notifications, Personalization, Natural Language Generation, Recommender Systems, Causal Inference, Retrieval-Augmented Generation, Reward Modeling}

\maketitle

\section{Introduction}

In the contemporary digital economy, user attention is scarce, interruptible, and intensely contested. Consumer platforms invest heavily in acquisition, but long-term value depends just as much on retention, reactivation, and repeated high-value engagement. Push notifications have therefore become a central operational mechanism for re-engaging users outside the application surface. On large consumer platforms, daily notification volumes range from millions to billions, and small changes in notification quality can materially affect click-through rates (CTR), conversion outcomes, and the broader health of the platform ecosystem \cite{bie2026pushgen} [T1]. Poorly targeted or repetitive notifications do not merely underperform; they actively erode trust, accelerate opt-outs, and contribute to long-term channel degradation \cite{ohly2023effects} [T1].

Historically, platforms have relied on two broad strategies for notification composition: manual copywriting and template-based generation. Most production notification systems today remain template-driven. Even when personalization is present, it is often limited to slot filling, such as inserting a user’s preferred category, recent item, or location into a fixed message pattern. This makes notifications relevant in a narrow sense, but rarely intelligent in how they communicate. 

The central observation of this paper is that modern notification stacks have invested heavily in three of four key decisions (audience selection, content ranking, send-time optimization), but the fourth decision, \textit{how to communicate the recommendation}, remains the least optimized stage. Consider a user with a historical preference for pizza receiving a dinner-time notification at 6 PM. A conventional system may send a static message such as ``Check out Tony's Pizza. Order now.'' This leaves engagement value on the table. An LLM-enabled system, by contrast, can incorporate temporal context, prior taste signals, and controlled novelty to generate a richer message such as: ``Perfect rainy lunch day---Tony's Margherita is 25 min away. Your usual, rated 4.8$\star$.'' In this setting, the value of the LLM is not in identifying the user or the category alone, but in producing a more intelligent, context-sensitive, and action-inducing framing of the recommendation.

Yet enthusiasm for LLM-based messaging has also created analytical confusion. In many production pipelines, the notification message is only one stage in a larger decision stack. Reported gains in engagement are often attributed broadly to ``AI notifications'' or ``LLM personalization,'' even when the observed uplift may be driven primarily by upstream improvements in targeting or timing. This conflation obscures both scientific understanding and engineering decision-making. 

This paper argues that message quality is an independent, underinvested lever and that LLMs provide a principled mechanism for improving it. The aim is not merely to summarize examples, but to clarify where language generation fits within the broader notification stack, what technical patterns currently dominate practice, how those patterns should be evaluated, and under what conditions LLM generation is actually worth deploying.

\subsection{Contributions}

This paper makes three contributions absent from prior work:

\begin{enumerate}
    \item \textbf{Message Quality Framework with Empirical Grounding:} A six-dimension definition of notification message quality (contextual relevance, clarity, actionability, novelty handling, linguistic freshness, persuasive appropriateness) with an evidence-backed account of how LLMs improve each dimension.
    \item \textbf{Architectural Attribution:} A systematic disentanglement of the message generation layer from adjacent pipeline components (targeting, ranking, timing), specifying which observed gains can and cannot be attributed to language quality.
    \item \textbf{Binding-Constraint Framework:} A three-criterion decision framework specifying when LLM generation is the binding constraint and when templates remain superior, including a principled argument for \textit{when not to use LLMs at all}.
\end{enumerate}

\section{Why Message Quality Is an Independent Optimization Target}

\subsection{The Structural Ceiling of Template-Based Messaging}

Template-based systems (e.g., ``[User], your [Item] is waiting!'') hit a structural ceiling because they cannot adapt their rhetorical strategy to the context of the user. The template assumes that the fact of relevance is sufficient to drive action. In reality, users respond to \textit{how} relevance is framed. 

\subsection{A Six-Dimension Definition of Message Quality}

We propose that notification message quality is not a scalar property (like CTR) but a vector of six dimensions. Table \ref{tab:quality_framework} contrasts the capabilities of template-based systems versus LLM-based generation across these dimensions.

\begin{table*}[t]
\caption{The Six-Dimension Message Quality Framework}
\label{tab:quality_framework}
\begin{tabular}{p{3.5cm}p{5.5cm}p{5.5cm}}
\toprule
\textbf{Dimension} & \textbf{Template Capability} & \textbf{LLM Capability} \\
\midrule
\textbf{Contextual Relevance} & Weak---only static slot values & Strong---composes multiple signals naturally \\
\textbf{Clarity} & Variable---slot grammar often awkward & Flexible---optimizes phrasing for brevity \\
\textbf{Actionability} & Formulaic call-to-action & Context-sensitive framing of action \\
\textbf{Novelty Handling} & Poor---novelty reads as irrelevance & Better---bridges from known to adjacent \\
\textbf{Linguistic Freshness} & Low---same structure across exposures & High---semantic variety across exposures \\
\textbf{Persuasive Appropriateness} & Neutral---limited expressive range & Variable---requires explicit guardrails \\
\bottomrule
\end{tabular}
\end{table*}

\subsection{How LLMs Improve Each Quality Dimension: Mechanisms and Evidence}

\subsubsection{Contextual Relevance}
LLMs can synthesize disparate context signals (weather, time, past orders) into a coherent sentence. DoorDash reported a +1.8\% engagement lift from LLM framing alone, holding the user, item, and timing constant \cite{sinha2024doordash} [T2]. The PushGen system demonstrated 3.2$\times$ greater lexical diversity than baselines while maintaining relevance \cite{bie2026pushgen} [T1].

\subsubsection{Clarity}
Because LLMs can rewrite rather than just insert, they can optimize for brevity. A recent e-commerce deployment found that LLM-generated messages were 15\% shorter on average while achieving +12.5\% CTR and +8.3\% CVR \cite{acm2025ecommerce} [T1].

\subsubsection{Actionability}
Templates rely on formulaic calls-to-action (``Complete your order''). LLMs can frame the action contextually: ``Still thinking about those running shoes? They're in stock and you're \$12 from free shipping.''

\subsubsection{Novelty Handling}
When introducing a novel item, templates often fail because the user has no prior schema for it. LLMs can bridge the gap. The PushGen system's +14.08\% CTR improvement included explicit novelty injection \cite{bie2026pushgen} [T1]. Instagram's diversity penalty improved CTR even while reducing overall notification volume by preventing novelty fatigue \cite{sun2025ranking} [T2].

\subsubsection{Linguistic Freshness}
Template CTR typically decays 15--25\% over 30 days due to ad fatigue. LLM messages in e-commerce maintained 18\% higher engagement after 4 weeks by varying the semantic structure across exposures \cite{acm2025ecommerce} [T1].

\subsubsection{Persuasive Appropriateness}
Dynamic persuasive frame selection (e.g., switching between social proof, scarcity, and utility) yielded +12.5\% CTR and +8.3\% CVR, with gains concentrated among the ``movable middle'' users \cite{acm2025ecommerce} [T1].

\subsection{Aggregate Benefit: The Compounding Effect}

These six dimensions compound. The +12--14\% gains observed in the most successful deployments exceed what any single dimension predicts, suggesting that freshness, relevance, and actionability interact multiplicatively.

\subsection{When LLM Deployment Is and Is Not Warranted}

LLMs are not universally required. Deployment is warranted under three conditions: (1) the content admits multiple framings, (2) user response is sensitive to linguistic nuance, and (3) there is sufficient context for grounded personalization. Table \ref{tab:evidence} summarizes the production evidence mapped to these dimensions.

\begin{table*}[t]
\caption{Summary of Production Evidence for LLM-Based Notifications}
\label{tab:evidence}
\begin{tabular}{p{3cm}p{2.5cm}p{5cm}p{4cm}c}
\toprule
\textbf{System} & \textbf{Domain} & \textbf{Primary Quality Dimensions} & \textbf{Key Result} & \textbf{Tier} \\
\midrule
\textbf{PushGen \cite{bie2026pushgen}} & Social Media & Freshness, Novelty Handling & +14.08\% CTR & [T1] \\
\textbf{Instagram \cite{sun2025ranking}} & Social Media & Freshness, Novelty Handling & Better CTR, lower volume & [T2] \\
\textbf{E-Commerce \cite{acm2025ecommerce}} & Retail & Actionability, Persuasive Framing & +12.5\% CTR, +8.3\% CVR & [T1] \\
\textbf{DoorDash \cite{sinha2024doordash}} & Food Delivery & Contextual Relevance, Clarity & +1.8\% engagement & [T2] \\
\bottomrule
\end{tabular}
\end{table*}

\section{Review Methodology}

To support a structured and transparent synthesis, this survey followed a systematic review methodology. We identified 142 records; excluded 87 (off-topic/pre-2018/non-English); assessed 55 full texts; excluded 27 for insufficient detail; and retained 28 primary sources (18 T1, 7 T2, 3 T3).

\subsection{Search Strategy and PRISMA Flow}

We conducted targeted searches across arXiv, ACM Digital Library, IEEE Xplore, and ScienceDirect, along with engineering publications. The review covered January 2018 through March 2026. The screening process followed a formal PRISMA flow (Figure \ref{fig:prisma}).

\begin{figure}[h]
    \centering
    \includegraphics[width=\linewidth]{../figures/fig1_prisma.png}
    \caption{PRISMA flow diagram detailing the literature search, screening, and inclusion process.}
    \label{fig:prisma}
\end{figure}

\subsection{Evidence Tiering}

Sources were organized into three evidence tiers: [T1] Peer-reviewed papers; [T2] Official engineering blogs; [T3] Contextual vendor documentation.

\section{Disentangling LLM Contributions from Adjacent Systems}

A recurring analytical error is treating notification performance as if it were primarily a function of the copy itself. Table \ref{tab:attribution} clarifies the architectural attribution.

\begin{table}[h]
\caption{Architectural Attribution Matrix}
\label{tab:attribution}
\begin{tabular}{p{3cm}p{4.5cm}}
\toprule
\textbf{System Component} & \textbf{LLM Relevance} \\
\midrule
\textbf{Audience Selection} & None (a targeting problem) \\
\textbf{Content Ranking} & None (a relevance problem) \\
\textbf{Send-Time Opt.} & None (a timing problem) \\
\textbf{Message Generation} & Direct: the message quality problem \\
\textbf{Candidate Ranking} & Direct: reward model selects best \\
\bottomrule
\end{tabular}
\end{table}

LLMs should be evaluated against message-layer baselines, not against overall system performance. Three deployments offer clean attribution: DoorDash (held item constant, varied framing only $\rightarrow$ +1.8\%); PushGen (pairwise reward model controlling for item/user effects); and E-Commerce LLM (within-user A/B, same user received LLM vs template for same categories over time).

\section{Core Architectural Frameworks and Technical Tradeoffs}

\begin{table*}[t]
\caption{Core Architectural Frameworks and Quality Dimensions Served}
\label{tab:synthesis}
\begin{tabular}{p{3.5cm}p{4.5cm}p{4cm}p{3cm}}
\toprule
\textbf{Architectural Pattern} & \textbf{Quality Dims Served} & \textbf{Main Weakness} & \textbf{Key Failure Mode} \\
\midrule
\textbf{Template / Slot Filling} & (Baseline) & Low diversity & Message fatigue \\
\textbf{RAG-Grounded LLM} & Relevance, Clarity & Retrieval latency & Faithfulness hallucination \\
\textbf{PEFT / LoRA} & Freshness, Persuasion & Requires style curation & Overfitting \\
\textbf{Pairwise Reward} & Actionability, Novelty & Exposure bias & Reward hacking \\
\textbf{Diversity Penalty} & Freshness, Novelty & Can suppress CTR & Under-delivery \\
\textbf{Contextual Bandit} & (Timing, adjacent) & Exploration cost & Wrong-time delivery \\
\textbf{Budget Router} & (Routing, meta) & CLV estimation dependency & Over-serving \\
\bottomrule
\end{tabular}
\end{table*}

\subsection{Retrieval-Augmented Generation (RAG)}
RAG controls \textit{what} the model knows. It is a natural fit for notifications because they are highly contextual \cite{lewis2020retrieval} [T1]. However, retrieval granularity affects both relevance and latency \cite{mixofgranularity2025} [T1]. A critical risk is faithfulness hallucination (e.g., fabricated discount percentages, incorrect ETAs) \cite{benchmarking2025} [T1]. Factuality guards are architecturally essential.

\subsection{Parameter-Efficient Fine-Tuning (PEFT/LoRA)}
PEFT controls \textit{how} the model speaks. A single platform may need dozens of voices. LoRA allows $W_{merged} = W_0 + BA$, providing zero inference latency overhead \cite{hu2022lora, han2024parameter} [T1]. Tone becomes a continuous, learnable parameter.

\subsection{Pairwise Reward Modeling}
PushGen generates 4--8 candidates per event and ranks them via $P(m_i > m_j | u, c)$ \cite{bie2026pushgen} [T1]. This ``generate-then-rank'' paradigm decouples creativity from quality control. However, it is vulnerable to exposure bias \cite{joachims2018deep} [T1].

\subsection{Send-Time Optimization (STO)}
Timing is a first-class pipeline component, not a post-hoc parameter. Duolingo models STO as a ``sleeping recovering bandit'' using Thompson Sampling \cite{yancey2020sleeping} [T1].

\section{Domain-Specific Applications}

The meaning of a ``good'' notification varies by industry. Table \ref{tab:domain} maps these differences.

\begin{table*}[t]
\caption{Domain-Specific Mapping}
\label{tab:domain}
\begin{tabular}{p{2.5cm}p{4cm}p{4cm}p{4cm}}
\toprule
\textbf{Domain} & \textbf{Primary Quality Dims} & \textbf{LLM Primary Value} & \textbf{Dominant Risk} \\
\midrule
\textbf{Social Media} & Freshness, Novelty Handling & Semantic variety, reduced fatigue & Over-frequency, tone mismatch \\
\textbf{Food Delivery} & Contextual Relevance, Clarity & RAG-grounded temporal framing & False urgency, stale ETA claims \\
\textbf{E-Commerce} & Actionability, Persuasion & Dynamic frame selection & Dark patterns, hallucinated prices \\
\bottomrule
\end{tabular}
\end{table*}

\section{When LLM Text Generation Is Not the Binding Constraint}

\subsection{When Templates Remain Superior}
For transactional or low-variance alerts (e.g., ``Your order has shipped''), contextual relevance is already high, clarity is maximized by a single fact, novelty is irrelevant, freshness is immaterial, and persuasion is best served by constrained expression.

\subsection{The Three-Criterion Decision Framework}
LLM generation is appropriate only when \textbf{all three} criteria are met:
\begin{enumerate}
    \item \textbf{Framing Variance:} Does the content admit multiple plausible, meaningfully different framings?
    \item \textbf{Linguistic Sensitivity:} Is user response sensitive to \textit{how} (not just \textit{whether}) the recommendation is framed?
    \item \textbf{Context Richness:} Does the system have sufficient grounded context for non-trivial composition?
\end{enumerate}

\subsection{The Marginal Value Question}
The variance in reported lifts (+1\% to +14.5\%) is driven primarily by baseline sophistication. PushGen (+14.08\%) and e-commerce (+12.5\%) replaced underdeveloped baselines; DoorDash (+1.8\%) competed against substantial upstream optimization.

\section{Proposed Unified Architectural Framework}

Figure \ref{fig:pipeline} illustrates the proposed unified LLM-based notification pipeline.

\begin{figure*}[t]
    \centering
    \includegraphics[width=\textwidth]{../figures/fig2_pipeline.png}
    \caption{Unified LLM-based notification pipeline. The Budget-Aware Router uses the three-criterion framework to decide per-event whether to invoke the LLM path or fall back to templates.}
    \label{fig:pipeline}
\end{figure*}

\subsection{Architectural Flow}
(1) RAG retrieves context; (2) generator produces 4--8 candidates via LoRA; (3) Factuality and Policy Guard filters them; (4) Pairwise Reward Model ranks them; (5) Diversity and Frequency Control prevents fatigue; (6) STO Bandit optimizes timing; (7) Online Learning Layer updates the models.

\subsection{Failure Modes}
If the guardrail is too strict, it throttles throughput; if too loose, it allows factual errors. If the reward model has exposure bias, it selects the same ``safe'' candidates. If the Budget Router over-routes, it violates latency SLAs; if it under-routes, high-value users receive templates.

\section{Evaluation Frameworks Beyond Short-Term CTR}

\subsection{Identification Problem}
SUTVA violations are common in networked environments \cite{luo2024survey} [T1]. We must distinguish three treatment layers: item selection, timing, and message realization.

\subsection{Offline Evaluation Is Structurally Biased}
Novel prompts have no counterfactual labels, leading to selection bias at the population level.

\subsection{Heterogeneous Treatment Effects}
Tu et al. (2021) showed effects vary drastically across cohorts \cite{tu2021personalized} [T1].

\subsection{Multi-Objective Evaluation}
Evaluation requires three layers: (1) Primary (CTR, CVR, order value); (2) Guardrail (opt-out rate, mute actions, dismissals); (3) Long-horizon (7/30-day retention, fatigue trajectories). Guardrail violations must be treated as hard constraints.

\section{Ethical Frameworks and Governance}

\subsection{Persuasive Asymmetry}
Susser et al. (2019) define online manipulation as bypassing rational agency \cite{susser2019technology} [T1]. False scarcity is manipulation.

\subsection{Fairness and Differential Susceptibility}
Reward models must not systematically exploit the vulnerabilities of specific demographic groups \cite{wang2023survey} [T1].

\subsection{Accountability}
Under DSA Article 40 and GDPR Article 22, audit logs and factuality checks are part of core production design \cite{fabbri2023ethical} [T1].

\section{Threats to Validity}

\begin{enumerate}
    \item \textbf{Publication Bias:} Industry labs are incentivized to publish successful A/B tests and suppress negative results.
    \item \textbf{Temporal Degradation:} Rapid LLM evolution means 2024 constraints may be resolved by 2026.
    \item \textbf{Domain Representation:} Heavy focus on Social Media and E-Commerce.
    \item \textbf{Reproducibility:} Reliance on proprietary data and lack of open benchmarks.
    \item \textbf{Framework Limitations:} Pipeline trades latency/compute for quality; performance depends heavily on the Budget-Aware Router's accuracy.
\end{enumerate}

\section{Conclusion}

LLMs improve notification systems by improving \textit{message quality}, not by solving the targeting, ranking, or timing problems. CTR improvements range from +1\% to +14.5\% depending on baseline maturity. The binding-constraint framework offers practitioners a principled basis for making deployment decisions, and for deciding \textit{when not to deploy LLMs at all}.

\bibliographystyle{ACM-Reference-Format}
\bibliography{references}

\end{document}
